\documentclass[aspectratio=169,10pt]{beamer}
\usepackage[brazil]{babel}

\usetheme{UELevento}

\title[Introdução ao LaTeX]{Introdução ao LaTeX}
\subtitle{Documentos científicos sem dor de cabeça}
\evento{Semana da Computação 2026}
\datalocal{12 de março de 2026 \textbullet{} Auditório do CCE, UEL}
\palestrante[foto-exemplo.png]{Prof.\ Dr.\ Nome do Palestrante}{Departamento de Computação, UEL}
\palestrante{Nome da Palestrante}{Universidade Parceira}
\apoio{\textbf{Departamento de Computação} \textbullet{} Centro de Ciências Exatas \textbullet{} UEL}

\begin{document}

\maketitle

\begin{frame}{Programação}
  \begin{itemize}
    \item 19h00 \textemdash{} Abertura
    \item 19h15 \textemdash{} Palestra: primeiros passos
    \item 20h15 \textemdash{} Perguntas e conversa
  \end{itemize}
\end{frame}

\begin{frame}{Por que usar LaTeX?}{Três razões}
  \begin{itemize}
    \item Fórmulas e referências bem formatadas
    \item O texto é versionável
    \item Modelos prontos para a universidade
  \end{itemize}
  \begin{block}{Para levar}
    Comece pelo modelo da sua instituição e mude só o conteúdo.
  \end{block}
\end{frame}

\begin{frame}{Exemplo}
  \[ \int_0^1 x^2\,dx = \frac{1}{3} \]
\end{frame}

\obrigado[nome.palestrante@exemplo.uel.br]

\end{document}
